\chapter{改一个标题的字号，后面跟着全乱了}

\begin{notebox}
本章解决 Word 时代的第一笔血债：格式失控。学完这章，你改格式永远只改一行。
\end{notebox}

\painhead{烦恼}

“标题的字号又变了，全文得一个一个改。”

\painhead{现象 · 这些场景你肯定见过}

\textbf{场景一}：你的毕业论文写到第三章，导师发话：“所有一级标题换成黑体四号。”你在 Word 里选中“第一章 绪论”，改字号，改字体；然后是第二章；然后你发现第三章忘改了；然后你发现还有个“3.1”算几级标题来着？

\textbf{场景二}：手动编号灾难。你给标题敲了“一、二、三”，后来往中间插了一章，于是“三”要改“四”，“四”要改“五”……全文的编号一个一个挪，挪完发现交叉引用的编号又对不上了。

\textbf{场景三}：室友拿你的模板写他的部分，交回来一看——标题字号三种、编号风格两套（有“1”的有“一”的）、行距各不相同。合并文档的下午，你贡献给了格式刷。

这三个场景共享同一个病根：\textbf{在 Word 里，内容和格式长在一起}。你看到的每个字都自带格式，改格式就得找到每一处去改。

LaTeX 的解法一句话：\textbf{内容和格式分开。你只声明“这是个标题”，至于标题长什么样，文档开头一句话统一规定。}

\painhead{解决线路}

\subsection*{2.1 最小例子：声明式写文档}

新建文件，敲进去，XeLaTeX 编译：

\begin{lstlisting}
\documentclass[UTF8]{ctexart}
\begin{document}

\section{绪论}

本章介绍研究背景。

\subsection{研究背景}

近年来，大家越来越关注这个问题。

\subsubsection{国外研究现状}

外国学者也研究过。

\section{方法}

本章介绍研究方法。

\end{document}
\end{lstlisting}

编译出来，你会看到：\textbf{所有标题自动编号}（1、1.1、1.1.1、2……），字号字体层级分明，还各有合适的间距——你没做任何设置，这些都是 LaTeX 的默认排版，而且说实话，比大多数同学手动调的都专业。

\textbf{逐行拆解}：

\begin{itemize}
\setlength{\itemsep}{2pt}
\item \texttt{\textbackslash{}section\{绪论\}}：声明“这里是个一级标题，名字叫绪论”。注意你\textbf{没有}说“绪论这两个字用黑体三号”——格式的事，LaTeX 自己定；
\item \texttt{\textbackslash{}subsection\{研究背景\}}：二级标题，自动编成 1.1；
\item \texttt{\textbackslash{}subsubsection\{...\}}：三级标题；
\item 最妙的是：\textbf{你现在往中间再插一节，后面的编号全自动重排}。场景二的编号灾难，从根源上不存在了。
\end{itemize}

\subsection*{2.2 改字号：改一个词，全文生效}

导师说“正文用小四”。Word 里你要全选正文再设置（还容易把标题也选进去）。LaTeX 里，改第一行：

\begin{lstlisting}
\documentclass[UTF8, 12pt]{ctexart}
\end{lstlisting}

\texttt{[...]} 里的东西叫\textbf{文档类选项}，就是给 LaTeX 的统一设置。\texttt{12pt} 表示正文用 12 磅字号，所有标题会按比例自动缩放，依然层级分明。

再来几个常用选项，直接填进那个方括号，逗号隔开：

\begin{table}[htbp]
\centering
\small
\begin{tabularx}{\linewidth}{@{}*{2}{>{\raggedright\arraybackslash}X}@{}}
\toprule
\textbf{选项} & \textbf{作用} \\
\midrule
\texttt{UTF8} & 中文编码，写中文必带 \\
\texttt{12pt} & 正文 12 磅（默认 10pt；可选 10/11/12） \\
\texttt{a4paper} & A4 纸（默认就是，可显式写上图个安心） \\
\texttt{twoside} & 双面打印的左右边距交替 \\
\bottomrule
\end{tabularx}
\end{table}

\textbf{体会一下这个感觉}：全文几千字的字号，改了第一行的四个字符。这就是“格式集中管理”——你后面写的每一章，都自动继承文档开头定下的规矩。

\subsection*{2.3 长文档怎么办：report 和 book}

\texttt{article} 不分“章”，只分“节”。要写毕业论文那种带“第 X 章”的，换文档类：

\begin{lstlisting}
\documentclass[UTF8, 12pt]{ctexrep}
\end{lstlisting}

编译后你会发现标题变成“第 1 章 绪论”了，\texttt{\textbackslash{}section} 自动降级成章内小节（编成 1.1）。

三个文档类一张表：

\begin{table}[htbp]
\centering
\small
\begin{tabularx}{\linewidth}{@{}*{3}{>{\raggedright\arraybackslash}X}@{}}
\toprule
\textbf{文档类} & \textbf{适用} & \textbf{顶层结构} \\
\midrule
\texttt{ctexart} & 课程作业、实验报告、短文 & 1 节 \\
\texttt{ctexrep} & 毕业论文、学位论文 & 第 1 章 \\
\texttt{ctexbook} & 书、超长文档 & 第 1 章（另起页等书籍规范） \\
\bottomrule
\end{tabularx}
\end{table}

\textbf{选型的土办法}：有“章”就 rep/book，没有就 art。拿不准就 art。

\subsection*{2.4 例外情况：我就是想要个不编号的标题}

比如“参考文献”“致谢”这种不参与编号的标题，在命令后面加个 \texttt{*}：

\begin{lstlisting}
\section*{致谢}
\end{lstlisting}

带星号的版本不编号、也不进目录（想让它进目录有办法，第 3 章讲目录时说）。

\painhead{Word 做法 vs LaTeX 做法 · 对照小结}

\begin{table}[htbp]
\centering
\small
\begin{tabularx}{\linewidth}{@{}*{3}{>{\raggedright\arraybackslash}X}@{}}
\toprule
\textbf{事项} & \textbf{Word 做法} & \textbf{LaTeX 做法} \\
\midrule
改标题格式 & 逐个选中改，或调样式表（大部分人不会用样式） & 改文档类或选项，一处生效 \\
章节编号 & 手打，或用多级列表（配半天） & \texttt{\textbackslash{}section} 自动编号 \\
插入新章节 & 后面编号全部手动顺延 & 零操作，自动重排 \\
不同级别标题 & 自己记"一级黑体三号、二级黑体四号……" & 层级自带，格式自带 \\
队友格式不统一 & 合稿后格式刷刷到死 & 格式在文档类里定死了，想乱都乱不了 \\
\bottomrule
\end{tabularx}
\end{table}

\painhead{翻车合集}

\begin{table}[htbp]
\centering
\small
\begin{tabularx}{\linewidth}{@{}*{3}{>{\raggedright\arraybackslash}X}@{}}
\toprule
\textbf{症状} & \textbf{原因} & \textbf{解法} \\
\midrule
\texttt{! LaTeX Error: Lonely \textbackslash{}section} 或类似 & \texttt{\textbackslash{}section\{\}} 花括号忘了闭合 & 检查命令的大括号是否成对 \\
编出来的中文是空心方块或问号 & 忘了 \texttt{UTF8} 选项，或文件保存编码不对 & 文档类加 \texttt{UTF8}；TeXworks 默认 UTF-8 一般不用动 \\
标题编到了 4 级没反应 & \texttt{article} 最多支持到 \texttt{\textbackslash{}subsubsection}（4 级） & 更深的层级用列表，别硬降级 \\
想要"第 1 章"但出来了"1" & 文档类用了 \texttt{ctexart} & 换 \texttt{ctexrep} 或 \texttt{ctexbook} \\
两级标题之间空了一大片 & 正常现象，LaTeX 的默认间距 & 别手动塞空行去"修"，越修越乱 \\
\bottomrule
\end{tabularx}
\end{table}

\section*{练习}

\begin{enumerate}
\setlength{\itemsep}{2pt}
\item 用 \texttt{ctexart} 写一份三节的小报告：《我的一周》，每节下面至少一个小节。观察编号；
\item 在第一节和第二节之间\textbf{插入一个新节}，重新编译，确认后面编号全部自动顺延；
\item 把文档类改成 \texttt{ctexrep}，看看“节”变成“章”后的编号长什么样，然后改回来；
\item 给《我的一周》加一个不编号的“附：本周购物清单”标题（星号版本）；
\item 灵魂拷问：打开你上学期交的某份 Word 作业，检查标题格式是否真的全篇统一。答案很可能是“并没有”——记住这个痛。
\end{enumerate}

到这里，你已经能排出结构清晰、编号可靠、格式统一的文档骨架了。下一章收拾第二笔血债：\textbf{目录和页码要手动刷新}——顺便，你会发现“编译两遍”的真理。
